\chapter{Conclusion}
\label{ch:conclusion}

\section{Project Summary}
This project successfully developed an AI-based proof-of-concept system for the non-intrusive extraction of swimming kinematics from uncalibrated, handheld poolside video footage. Driven by the need to overcome severe visual distortions present in aquatic environments, such as refraction, turbulence, and occlusion, a modular computer vision pipeline was proposed. By introducing a novel dynamic pixel-to-meter scaling methodology utilizing custom-designed underwater markers, the system demonstrated the feasibility of extracting real-world metric data without relying on rigid camera setups. The integration of advanced deep learning foundation models, including SAM 2 for temporal segmentation, YOLO26 for marker detection, and ViTPose for human pose estimation, provided a robust framework for tracking swimmer movements and addressing domain-specific data scarcity.

\section{Achievement of Objectives}
The project successfully met its core objectives:
\begin{enumerate}
    \item \textbf{Marker Design:} A highly visible, modular 3D-printed marker system was developed to withstand underwater turbulence and provide a reliable calibration reference at 2.5\,m intervals.
    \item \textbf{Dataset Generation (MITL):} A Model-in-the-Loop annotation pipeline leveraging SAM 2, statistical outlier filtering (IQR), and human verification was implemented to construct a domain-specific dataset for training a custom YOLO26s detector.
    \item \textbf{Dynamic Scale Extraction:} A refraction-aware pipeline was established, utilizing Median Absolute Deviation (MAD) filtering to stabilize the pixel-to-meter scale derived from the known marker spacing.
    \item \textbf{Kinematic Metric Extraction:} Combining ego-motion compensation, ViTPose localization, and a Kalman filter smoothing stage, the system successfully produced absolute positional and velocity profiles from the raw tracking output. 
\end{enumerate}

\section{Limitations and Challenges}
Despite proving the feasibility of the concept, the evaluation highlighted several inherent challenges in the aquatic domain. The positional accuracy evaluation on a 50\,m breaststroke sequence yielded a Mean Absolute Error (MAE) of 1.005\,s and a Root Mean Square Error (RMSE) of 1.210\,s across 18 checkpoints. A primary limitation was the non-linear fluctuation in scale estimation on the outbound leg, leading to alternating periods of overestimated and underestimated swimming speeds. Furthermore, while macroscopic centroid tracking was successfully extracted, tracking distal keypoints such as the wrists remained noisy underwater and exhibited substantial high-frequency artefacts. Consequently, fine-grain kinematic noise limited the utility of the current system for precision technique analysis. Lastly, the evaluation relied on single-observer manual ground truth split times, which inherently introduces a limitation regarding inter-rater reliability.

\section{Future Work}
Building upon this proof of concept, several avenues for future research and refinement have been identified:
\begin{itemize}
    \item \textbf{Enhancement of Pose Estimation:} Future iterations should target the refinement of distal keypoint tracking (particularly for wrists) to enable detailed, centroid-relative technique analysis.
    \item \textbf{Advanced Scale Stabilisation:} Methods to mitigate the non-linear scaling fluctuations caused by severe refraction should be explored to improve the continuous monotonic accuracy of the estimated swimming speed. 
    \item \textbf{Broader Dataset Validation:} While the proof of concept successfully evaluated a breaststroke sequence, future work should validate the system's robustness across different strokes, variable lighting, and different pools.
    \item \textbf{Automated Ground Truth Collection:} Replacing manual, single-observer timestamp recording with sensor-based timing systems (such as touchpads) would provide a more objective baseline for positional accuracy evaluation.
\end{itemize}